% 答案与解析 · 第 1 章（对应 ../handout/chapters/01）
\section{集合与常用逻辑用语}

\begin{ansblock}{练习 1.1}
\begin{enumerate}
  \item （1）$\{1,2,3,4,5\}$；
  （2）$\{-1,1\}$；
  （3）$\{x\mid x>3\}$（或 $(3,+\infty)$）。
\end{enumerate}
\end{ansblock}

\begin{ansblock}{练习 1.2（含解析）}
\begin{enumerate}
  \item $A=\{x\mid 2<x<3\}$。由 $A\subseteq B$ 需 $B$ 的右端点 $a\ge 3$，
        故 $a$ 的取值范围是 $[3,+\infty)$。
        \quad\textit{易错点：把 $A\subseteq B$ 写成 $a>3$，漏掉端点。}
  \item $q$ 的解集为 $(0,2)$，与 $p$ 相同，故 $p$ 是 $q$ 的充要条件。
  \item 不能。集合元素满足互异性，$\{1,1,2\}=\{1,2\}$；
        若题目要求「由 1,2 组成的集合」，应写成 $\{1,2\}$。
\end{enumerate}
\end{ansblock}

\begin{ansblock}{补充：三种解法的比较}
第 1.2 节例 1 还可用配方法：$x^2-2x-3=(x-1)^2-4$，由 $(x-1)^2-4<0$ 得
$|x-1|<2$，即 $-1<x<3$。配方法的好处是不需因式分解，对十字相乘不熟练的学生更稳。
\end{ansblock}
